% Source: midterm1_v1.html
\documentclass[11pt]{article}
\usepackage[margin=1in]{geometry}
\usepackage{amsmath}
\usepackage{amssymb}
\usepackage{siunitx}
\usepackage{enumitem}

\title{Midterm 1 (v1) -- Fluids, Oscillations, Waves, Sound}
\author{}
\date{}

\begin{document}
\maketitle

\section*{Midterm 1 (v1) -- Fluids, Oscillations, Waves, Sound}

\begin{enumerate}[leftmargin=*,label=\textbf{Question \arabic*.}]
\item \hfill \textit{/ 10 pts}

On a far planet the local acceleration due to gravity is 13.6 m/s$^{2}$ and the atmospheric pressure is 100,241 Pa. Calculate the pressure 6.8 m below the surface of a fluid having density 1,040 kg/m$^{3}$.



\emph{p} = \rule{2cm}{0.4pt} Pa

\textit{Correct answer: } 196,420.2

\item \hfill \textit{/ 10 pts}

Suppose a hydraulic press has an output piston that has 3.3 times the cross-sectional area as the input piston. If the input force is 419 N, what is the output force in kN?



\emph{F} = \rule{2cm}{0.4pt} kN

\textit{Correct answer: } 1.3827

\item \hfill \textit{/ 10 pts}

A ball of volume 12 m$^{3}$ is placed in a container of fluid of density 813 kg/m$^{3}$. Calculate the density of the object, if it floats and 8.7 m$^{3}$ of it is submerged in the fluid.



𝜌 = \rule{2cm}{0.4pt} kg/m$^{3}$

\textit{Correct answer: } 589.425

\item \hfill \textit{/ 10 pts}

How many minutes does it take to fill a 37-gallon bathtub if the faucet opening has an area of 1.3 cm$^{2}$ and water flows out of it at 3.1 m/s? (1 gal = 3.785 L = 3.785 × 10$^{-3}$ m$^{3}$.)



\emph{t} = \rule{2cm}{0.4pt} min

\textit{Correct answer: } 5.7918

\item \hfill \textit{/ 10 pts}

A horizontal pipe of cross-sectional area 0.34 m$^{2}$ at height 10 m carries water at speed 2 m/s. The pipe turns upward and levels out again to horizontal at height 13 m, where the pipe narrows to 0.12 m$^{2}$. If the pressure is 204 kPa at the original height, what is the pressure in kPa at the final height? The density of water is 1000 kg/m$^{3}$.



\emph{p} = \rule{2cm}{0.4pt} kPa

\textit{Correct answer: } 160.5444

\item \hfill \textit{/ 10 pts}

Suppose a horizontal block-spring system has a spring constant 1,330 N/m and block of mass 3.6 kg. Calculate the frequency in Hz of its harmonic motion.

\textit{Correct answer: } 3.0591

\item \hfill \textit{/ 10 pts}

A simple harmonic oscillator consists of a block of mass 3.4 kg and spring with an unknown spring constant \emph{k} on a frictionless horizontal surface. The block is moved a distance 0.49 mfrom equilibrium and released from rest. If the block's maximum speed at the equilibrium point is 5.4 m/s, calculate the spring constant \emph{k} of the block in N/m.



\emph{k} = \rule{2cm}{0.4pt} N/m

\textit{Correct answer: } 412.9279

\item \hfill \textit{/ 10 pts}

A block of mass 1.3 kg is attached to a spring of force constant 56 N/m. At the instant the spring is in its equilibrium position the block is moving in the positive direction at 1.1 m/s. Find the oscillation amplitude in cm.



\emph{A }= \rule{2cm}{0.4pt} cm

\textit{Correct answer: } 16.7599

\item \hfill \textit{/ 10 pts}

Calculate the length of a simple pendulum on Earth if its angular frequency is 3.61 rad/s.



 \emph{L} = \rule{2cm}{0.4pt} m

\textit{Correct answer: } 0.752

\item \hfill \textit{/ 20 pts}

Suppose a wave is given by \emph{y}(\emph{x},\emph{t})=5 cos(0.2𝜋\emph{t - }0.3𝜋\emph{x} + 0.25𝜋) with all quantities in MKS units. Determine the following in MKS units:

\textbf{(a)} amplitude
 [ Select ]

;

\textbf{(b)} phase shift
 [ Select ]

;

\textbf{(c)} wave number
 [ Select ]

;

\textbf{(d)} angular frequency
 [ Select ]

;

\textbf{(e)} frequency
 [ Select ]

;

\textbf{(f)} period
 [ Select ]

;

\textbf{(g)} wavelength
 [ Select ]


\textbf{(h)} propagation speed
 [ Select ]

;

\textbf{(i) }transverse \textbf{speed} at \emph{x} = 0, \emph{t} = 0:
 [ Select ]

\textit{Answer 1:}
\begin{itemize}
  \item 2 m
  \item 6 m
  \item 2.5 m
  \item \textbf{[correct]} 5 m
  \item 3 m
\end{itemize}
\textit{Answer 2:}
\begin{itemize}
  \item 0.25 rad
  \item 0.628 rad
  \item \textbf{[correct]} 0.785 rad
  \item 0.125 rad
\end{itemize}
\textit{Answer 3:}
\begin{itemize}
  \item 0.628 rad/m
  \item \textbf{[correct]} 0.942 rad/m
  \item 0.785 rad/m
  \item 0.15 rad/m
\end{itemize}
\textit{Answer 4:}
\begin{itemize}
  \item 0.942 rad/s
  \item 0.314 rad/s
  \item \textbf{[correct]} 0.628 rad/s
  \item 0.2 rad/s
\end{itemize}
\textit{Answer 5:}
\begin{itemize}
  \item \textbf{[correct]} 0.1 Hz
  \item 0.125 Hz
  \item 0.15 Hz
  \item 5 Hz
\end{itemize}
\textit{Answer 6:}
\begin{itemize}
  \item 0.942/m
  \item 6.67 s
  \item 0.15 /m
  \item 5 s
  \item 20 s
  \item 0.125/m
  \item 0.628/m
  \item \textbf{[correct]} 10 s
  \item 0.628 s
\end{itemize}
\textit{Answer 7:}
\begin{itemize}
  \item \textbf{[correct]} 6.67 m
  \item 0.942 m
  \item 5 m
  \item 10 m
\end{itemize}
\textit{Answer 8:}
\begin{itemize}
  \item \textbf{[correct]} 0.667 m/s
  \item 3.14 m/s
  \item 2.35 m/s
  \item 1.50 m/s
\end{itemize}
\textit{Answer 9:}
\begin{itemize}
  \item 3.14 m/s
  \item 10 m/s
  \item 6.67 m/s
  \item \textbf{[correct]} 2.22 m/s
\end{itemize}

\item \hfill \textit{/ 10 pts}

A wire with fundamental frequency 351 Hz and length 0.37 m is under tension of 1,403 N. Calculate the wire's mass per unit length in \textbf{grams/meter}.



$\mu=$ \rule{2cm}{0.4pt} g/m (Note the units!)

\textit{Correct answer: } 20.796

\item \hfill \textit{/ 10 pts}

A mine explodes off shore north of Cuba, creating a sound wave under water. How long does it take for the sound to register on Coast Guard sonar, 333 km away?

Density and bulk modulus of sea water: 1025 kg/m$^{3}$ and 2.34 x 10$^{9}$ Pa.



\emph{t} = \rule{2cm}{0.4pt} s

\textit{Correct answer: } 220.3932

\item \hfill \textit{/ 10 pts}

A pneumatic drill creates a decibel level of 117 dB at a distance of 12 m. Calculate the \textbf{sound intensity level }in dBat a distance of 51 m.

\textit{Correct answer: } 104.4322

\item \hfill \textit{/ 10 pts}

A pipe open at both ends has a length of 1 m. If the temperature of the room is 12°C, calculate the frequency in Hz of the \emph{n} = 1 harmonic.

Take the speed of sound to be $v=(331\;\mathrm{m/s})\sqrt{\frac{T}{273\;\mathrm{K}}}$ .

\textit{Correct answer: } 169.0982

\item \hfill \textit{/ 10 pts}

A motorcyclist is traveling 29 m/s due west relative the ground. A fire truck is traveling on the same road, going east at 40 m/s and away from the motorcyclist, having already passed it going the opposite direction. If the fire truck siren sounds at 1,094 Hz, what frequency does the motorcyclist hear? Sound speed in air: 350 m/s.

\textit{Correct answer: } 900.4462

\end{enumerate}

\end{document}
